\section{Finite relative-trace certificates}\label{sec:finite}
The same construction will be used for integral and quadratic dimensions.
In the quadratic case put $K=\Q(\sqrt q)$, with $q>1$ squarefree, and let
$F=q$ for $q\equiv1\pmod4$, and $F=4q$ otherwise. In the rational case
put $K=\Q$ and $F=1$. A rational computation is not interpreted as a
quadratic field with two embeddings.

\subsection{A finite set of joint traces}
Let $P$ be a finite set of rational primes containing those dividing $F$,
and fix a height $h\ge3$. For a root of unity $\zeta$ with order supported
on $P$, let
\[
 \tau_j(\zeta)=\frac{\Tr_{L/K}(\zeta^j)}{[L:K]},\qquad 1\le j\le h,
\]
where $L$ is any common cyclotomic field containing $K$ and $\zeta$.
The value is independent of enlarging $L$.

\begin{lemma}\label{lem:trace-box}
All these joint traces are contained in an explicitly computable finite
set. It suffices to use orders dividing
\begin{equation}\label{eq:trace-box}
 B_h=\prod_{p\in P}p^{\max(v_p(F),1)+\lfloor\log_p h\rfloor}
\end{equation}
and to add the all-zero vector. For $K=\Q$, the trace of a primitive
$n$th root is $\mu(n)/\varphi(n)$. For quadratic $K$, if $\chi$ is the
primitive quadratic character of conductor $F$, then
\begin{equation}\label{eq:gauss}
 \tau_1(\zeta_n^a)=\frac{\mu(n)}{\varphi(n)}+
 \begin{cases}
 \displaystyle\frac{\chi(a)\mu(g)\chi(g)\sqrt F}{\varphi(n)},
       &n=Fg,\ (F,g)=1,\\[3pt]
 0,&\text{otherwise},
 \end{cases}
\end{equation}
for $(a,n)=1$. Higher coordinates are obtained by reducing
$\zeta_n^{aj}$ to its exact order. The same sign $\chi(a)$ is used in
every nonrational coordinate of a joint trace.
\end{lemma}
\begin{proof}
For $F\nmid n$, the intersection $K\cap\Q(\zeta_n)$ is $\Q$, since the
conductor of $K$ is $F$, so the normalized trace is the normalized
Ramanujan sum. For $F\mid n$, average over the subgroup on which $\chi$
is one. Its indicator is $(1+\chi)/2$. The untwisted sum is $\mu(n)$,
and the induced quadratic Gauss sum is
$\chi(a)\mu(g)\chi(g)\sqrt F$ when $(F,g)=1$ and zero otherwise.
This gives \eqref{eq:gauss}, including nonsquarefree $g$.
For a power $j$, set $r=(n,j)$ and apply the formula at order $n/r$,
with primitive exponent $aj/r$; this proves the simultaneous sign rule.

Suppose $v_p(n)>\max(v_p(F),1)+\lfloor\log_p h\rfloor$ for some $p$.
For every $j\le h$, its reduced order still has $p$-valuation strictly
greater than both $v_p(F)$ and one. Its Ramanujan sum vanishes, and
so does its induced Gauss sum. Hence every coordinate is zero.
All other orders divide $B_h$. Taking both quadratic signs may retain
extra tuples when an exponent restriction is redundant; this enlarges
the necessary set and cannot create a false exclusion.
\end{proof}

The bound $B_h$ is a bound on \emph{trace types}, not on the exponent of
an unknown category. The categorical Cauchy theorem
\cite[Theorem 3.9]{BNRW}, applied to the modular center, allows $P$ to
be the prime divisors of $|N_{K/\Q}(D)|$ (of $D$ when $K=\Q$).
Thus unbounded possible twist orders cause no omitted trace types.

\subsection{Gram columns and indicator bounds}
Fix a low matrix $L$ satisfying \eqref{eq:low}. Each remaining central
simple has a restriction vector $v=(0,v_1,\ldots,v_{r-1})$ with
nonnegative integer entries, and dimension $z_v=\sum_i v_i d_i$.
The finite set of permitted vectors contains every actual vector if we
require
\begin{equation}\label{eq:high}
 v_iv_j\le Q_{ij},\qquad Q-vv^{\mathsf T}\succeq0,\qquad
 z_v\text{ a }d\text{-number},\qquad D/z_v\text{ an algebraic integer}.
\end{equation}
Indeed, the residual Gram matrix is a sum of the rank-one matrices of
all remaining restriction vectors. The dimension conditions are the
modular-center consequences recalled in \cite[Proposition 2.1]{DP}.
Finiteness follows from $v_i\le\lfloor\sqrt{Q_{ii}}\rfloor$.
Positive semidefiniteness is checked by principal minors; for a singular
$Q$ one also imposes its rational kernel before using a nonsingular
principal block. All calculations are rational or integral.

For $j\ge3$ let
\[
 H_{i,j}=[X_i^j:1],\qquad
 \mathcal J=\{3\le j\le h:K=\Q\text{ or }F\nmid j\}.
\]
The indicator is the trace of the order-$j$ rotation on an
$H_{i,j}$-dimensional space. Its normalized relative trace to $K$ is
rational for $j\in\mathcal J$, because
$K\cap\Q(\zeta_j)=\Q$. Denote it by $\bar\nu_j(X_i)$.
If $p(j)$ is the least prime divisor of $j$, then
\begin{equation}\label{eq:indicator-bound}
 -\frac{H_{i,j}}{p(j)-1}\le \bar\nu_j(X_i)\le H_{i,j}.
\end{equation}
To see this, each normalized trace of a rotation eigenvalue is
$\mu(n)/\varphi(n)$ for some $n\mid j$. If negative, its denominator is
at least $p(j)-1$. Summing the $H_{i,j}$ terms proves the bounds.
No numerical approximation to an indicator is involved.

For every permitted vector $v$ and joint trace type $\tau$ introduce a
nonnegative real variable $x_{v,\tau}$. It relaxes the integer number of
central simples having that data. The necessary equations are
\begin{align}
 \sum_{v,\tau}x_{v,\tau}v_iv_j&=Q_{ij},\label{eq:system-gram}\\
 \ell_i+\sum_{v,\tau}x_{v,\tau}v_i z_v\tau_1&=0,
                                                       \label{eq:system-one}\\
 \ell_i+\sum_{v,\tau}x_{v,\tau}v_i z_v\tau_2&=D\epsilon_i.
                                                       \label{eq:system-two}
\end{align}
Here $\epsilon_i$ runs over every second-indicator choice allowed in
Section~\ref{sec:inputs}. For $j\in\mathcal J$, also require
\begin{equation}\label{eq:system-higher}
 \ell_i+\sum_{v,\tau}x_{v,\tau}v_i z_v\tau_j=D u_{i,j},
 \qquad u_{i,j}\in\Q,\quad
 -H_{i,j}/(p(j)-1)\le u_{i,j}\le H_{i,j}.
\end{equation}
The rationality condition is a linear equation over $\Q$: writing
$D=D_0+D_1\sqrt q$ and the left side as $w_0+w_1\sqrt q$, it is
$D_1w_0-D_0w_1=0$. The bounds can then be imposed on $w_0/D_0$,
where $D_0>0$. For $K=\Q$ the rationality equation is void.
Expanding all field coefficients, adding nonnegative slack variables,
and clearing denominators gives an integer system
\begin{equation}\label{eq:farkas}
 A x=b,\qquad x\ge0.
\end{equation}
Zero rationality rows may be retained in the rational implementation.

\begin{theorem}[Finite-indicator criterion]\label{thm:finite}
Let $R$ be a commutative fusion ring whose FP dimension field is rational
or real quadratic. Enumerate all low matrices, all second-indicator
choices and all high columns as above. If each low matrix either fails
a condition in Section~\ref{sec:inputs}, or, for every remaining indicator
choice, its system \eqref{eq:farkas} has an integer vector $y$ satisfying
\begin{equation}\label{eq:dual-ray}
 A^{\mathsf T}y\ge0,\qquad b^{\mathsf T}y<0,
\end{equation}
then $R$ has no pseudo-unitary categorification.
\end{theorem}
\begin{proof}
A hypothetical categorification supplies one of the low matrices and one
of the indicator choices. Its remaining restriction vectors satisfy
\eqref{eq:high}; Lemma~\ref{lem:trace-box} contains all their joint traces.
Equations \eqref{eq:system-gram}--\eqref{eq:system-higher} follow from
adjunction and \eqref{eq:NS}, and the indicator bounds were just proved.
Consequently it supplies a nonnegative solution $x$. But
$0\le x^{\mathsf T}A^{\mathsf T}y=b^{\mathsf T}y<0$, a contradiction.
\end{proof}

All applications use height six. Retaining higher rational-order bounds
contains the third-indicator inequality rather than introducing another
criterion. The two-moment inequalities \eqref{eq:low-moment} are kept as
cheap initial consequences. A floating-point linear program may propose
a dual ray; only the cleared integer vector and the exact inequalities
\eqref{eq:dual-ray} are accepted. Failure to find such a ray is not a
positive categorification result.
